\section{Preliminary}
\label{sec:preliminary}

\subsection{Alpha Mining}

An alpha factor is an executable mapping from historical market
information to a cross-sectional signal:

\[
f:X_{t-\tau+1:t}\rightarrow v_t\in\mathbb{R}^{N}.
\]

Beyond its implementation, a factor reflects an underlying trading
mechanism that describes a market phenomenon, its applicable conditions,
and the expected direction of its effect. We view alpha discovery as
the process of identifying such mechanisms and realizing them as
executable factors.

Formally, let \(z\in\mathcal{Z}\) denote a point in a trading-semantic
space, where each \(z\) represents a structured description of a
candidate trading mechanism. A realization function
\[
R:\mathcal{Z}\rightarrow\mathcal{F}
\]
maps semantic descriptions into executable factors. Alpha mining then
aims to discover semantic candidates whose realizations achieve strong
factor quality:

\[
z^*=\arg\max_{z\in\mathcal{Z}}\mathcal{J}(R(z)).
\]

This formulation separates the semantic object being explored from the
implementation used for evaluation.

\subsection{Structured Semantic Plans}

\method{} instead searches over structured semantic plans, which describe a
candidate trading mechanism before code generation. A plan is a compact
specification with five fields:
\begin{itemize}[leftmargin=12pt, nosep]
  \item \textbf{Event}: the market phenomenon, such as breakout, volume
        expansion, or volatility compression;
  \item \textbf{Context}: the reference state in which the event is
        interpreted, such as a recent extreme, VWAP region, or cross-sectional
        quantile;
  \item \textbf{Qualities}: confirmation or filtering conditions, such as
        volume confirmation, multi-horizon consistency, or outlier removal;
  \item \textbf{Direction}: the assumed implication, such as continuation,
        reversal, or range oscillation;
  \item \textbf{Output}: the signal emission form, such as a continuous
        score, event decay, or cross-sectional rank.
\end{itemize}

Thus, a plan specifies \emph{what} event is observed, \emph{where} it is
conditioned, \emph{which} confirmations are required, \emph{what} direction is
implied, and \emph{how} the signal is emitted. The formal notation for this
schema space is deferred to Appendix~\ref{sec:appendix_schema_notation}. 